% Verossimilhança do GARCH(1,1): num dia (o de maior resíduo da janela) e em
% todos os dias. O GARCH escolhe omega, alfa e beta que tornam a soma dos
% log f(resíduo | sigma_t) a maior possível. Números em valores.tex.
\documentclass[border=4pt]{standalone}
% Preâmbulo comum dos gráficos que explicam os métodos (30/10/2026).
% Dados e números: dados/resultados/metodos_ilustrados_2026-10-30/, gerados por
% ferramentas/MetodosIlustrados.mq5 com as funções do próprio robô. Nenhum
% número é digitado aqui: os do texto vêm de valores.tex (\vX para o texto,
% \pX para coordenada).
\usepackage[T1]{fontenc}
\usepackage{amsmath}
\usepackage{mathptmx}  % Times, como as figuras do TG
\usepackage{pgfplots}
\usepgfplotslibrary{groupplots,dateplot,colormaps,fillbetween}
\usetikzlibrary{arrows.meta,positioning}
\pgfplotsset{compat=1.18}
\definecolor{cmod}{RGB}{0,90,181}     % modelo / ajuste
\definecolor{cdes}{RGB}{220,50,32}    % destaque, limite, valor crítico
\definecolor{cdad}{RGB}{120,120,120}  % dado observado
\newcommand{\arq}[1]{../../dados/resultados/metodos_ilustrados_2026-10-30/#1}
\input{\arq{valores.tex}}
% Figuras 08 a 14 (o robô operando): traço por barra da refação de 05/10/2026,
% resumido por ferramentas/scripts/robo_em_acao.sh. Macros \vR e \pR.
\newcommand{\arqr}[1]{../../dados/resultados/robo_em_acao_2026-10-30/#1}
\input{\arqr{valores.tex}}
\definecolor{cari}{RGB}{230,140,0}   % ARIMA puro
\definecolor{cok}{RGB}{0,140,70}     % compra liberada
\usetikzlibrary{shapes.misc}
\newcommand{\passo}[1]{\tikz[baseline=(n.base)]\node[circle,fill=black!75,text=white,inner sep=1pt,font=\footnotesize\bfseries](n){#1};}
% eixo de datas comum às figuras do traço (jul/2019 a jun/2021)
\pgfplotsset{traco/.style={date coordinates in=x, xmin=2019-07-01, xmax=2021-06-30,
    xtick={2019-07-01,2020-01-01,2020-07-01,2021-01-01,2021-06-30},
    xticklabels={jul/2019,jan/2020,jul/2020,jan/2021,jun/2021}}}
\pgfplotsset{
  base/.style={table/col sep=comma, grid=major, grid style={gray!20},
    tick label style={font=\small}, label style={font=\small},
    title style={font=\small\bfseries, align=center},
    every axis plot/.append style={line width=0.9pt},
    /pgf/number format/use comma, /pgf/number format/1000 sep={.}},
}
% faixa vertical nos pregões do traço em que a condição #1 vale, na cor #2,
% do piso #3 ao teto #4 do eixo
\newcommand{\faixa}[4]{%
  \addplot[name path=topo, draw=none, const plot, forget plot] table[x=data, y expr={(#1) ? #4 : #3}]{\arqr{dia_a_dia.csv}};
  \addplot[name path=piso, draw=none, forget plot] coordinates {(2019-07-01,#3) (2021-06-30,#3)};
  \addplot[#2, forget plot] fill between[of=topo and piso];}

\begin{document}
\begin{tikzpicture}
\begin{groupplot}[base, group style={group size=2 by 1, horizontal sep=1.7cm}, height=6.2cm]
\nextgroupplot[width=7.2cm, title={Um dia (\vDiaMax): com que $\sigma$\\o resíduo fica mais provável?},
  xlabel={Resíduo (\%)}, ylabel={$\log f$ (log da densidade normal)}, xmin=-10, xmax=10, ymin=-16, ymax=1,
  legend style={font=\footnotesize, at={(0.5,0.03)}, anchor=south, cells={anchor=west}}]
\addplot[cdad, dashed, line width=1.2pt] table[x=x, y expr=ln(\thisrow{amostral})]{\arq{gv_densidade.csv}};
\addplot[cmod, line width=1.4pt] table[x=x, y expr=ln(\thisrow{garch})]{\arq{gv_densidade.csv}};
\legend{$\sigma$ fixo $= \vGSigAm\%$, $\sigma_t$ do GARCH naquele dia $= \vSigT\%$}
\draw[cdes, line width=0.8pt, densely dotted] (axis cs:\pEpsMax,-16) -- (axis cs:\pEpsMax,1);
\addplot[only marks, mark=*, mark size=2.4pt, cdad] coordinates {(\pEpsMax,\pLogDensAm)};
\addplot[only marks, mark=*, mark size=2.4pt, cmod] coordinates {(\pEpsMax,\pLogDensT)};
\node[font=\footnotesize, anchor=west, text=cdad] at (axis cs:\pEpsMax,\pLogDensAm) {\ $\vLogDensAm$};
\node[font=\footnotesize, anchor=west, text=cmod] at (axis cs:\pEpsMax,\pLogDensT) {\ $\vLogDensT$};
\node[font=\footnotesize, anchor=north west, text=cdes, align=left] at (axis cs:\pEpsMax,0.6) {resíduo do dia: $\vEpsMax\%$};
\nextgroupplot[width=9.4cm, title={Todos os dias: $\sigma_t$ acompanha os períodos agitados\\$\log L = \sum_t \log f(\hat\varepsilon_t \mid \sigma_t) = \vGLogL$},
  date coordinates in=x, xticklabel={\month/\year}, ylabel={\%}, ymin=-8, ymax=8,
  xtick={2024-10-01,2025-04-01,2025-10-01,2026-04-01},
  legend style={font=\footnotesize, at={(0.5,-0.2)}, anchor=north, legend columns=2, draw=none, column sep=6pt}]
\addplot[name path=sup, cmod!60, line width=0.6pt, forget plot] table[x=data, y expr=2*\thisrow{sigma}]{\arq{gv_volatilidade.csv}};
\addplot[name path=inf, cmod!60, line width=0.6pt, forget plot] table[x=data, y expr=-2*\thisrow{sigma}]{\arq{gv_volatilidade.csv}};
\addplot[cmod!20, area legend] fill between[of=sup and inf];
\addplot[cdad, line width=0.5pt] table[x=data, y=residuo]{\arq{gv_volatilidade.csv}};
\legend{faixa $\pm 2\sigma_t$ do GARCH, resíduo do ARIMA $\hat\varepsilon_t$}
\end{groupplot}
\end{tikzpicture}
\end{document}
